%%
%% Capítulo 2: Desenvolvimento do trabalho
%%

\mychapter{Desenvolvimento}
\label{Cap:desenvolvimento}

Este capítulo descreve a revisão de literatura sobre os fundamentos teóricos utilizados, a metodologia adotada para a coleta, pré-processamento e modelagem dos dados, bem como os detalhes de implementação do pipeline experimental.

\section{Revisão de literatura}
\label{Sec:revisaoliteratura}

%% ================================================================
%% TODO: Seção sobre Criminologia Computacional e Predição de Crime
%% ================================================================
\subsection{TODO -- Predição espacial do crime}
\label{Sec:todo_crime_prediction}

% TODO: Esta subseção precisa ser escrita. Deve cobrir:
%
% 1. Análise espacial do crime e hotspots:
%    - Conceito de hotspot e concentração espacial do crime (já citado na
%      introdução via Weisburd 2015, mas precisa ser aprofundado aqui).
%    - Abordagens tradicionais: Kernel Density Estimation (KDE), modelos
%      de ponto (point-process), ARIMA espacial, etc.
%
% 2. Predição de crime como problema espaço-temporal:
%    - Representação em grids: como o espaço urbano é discretizado em
%      células de área fixa e o tempo em janelas regulares.
%    - Desafio da esparsidade: a maioria das células não apresenta crime
%      na maioria dos períodos, gerando dados altamente desbalanceados.
%    - Formulação do problema: entrada = sequência de mapas passados,
%      saída = mapa de risco futuro.
%
% 3. Aplicações de deep learning na predição de crime:
%    - Evolução dos métodos: de modelos estatísticos clássicos para
%      abordagens baseadas em redes neurais.
%    - Citar e discutir: Wang et al. (2019), Rotaru et al. (2022),
%      Albors Zumel et al. (2025) e outros trabalhos relevantes.
%    - Limitações reportadas na literatura: baixa precisão, vieses,
%      questões éticas (citar Papachristos 2022).
%
% Referências já na .bib que podem ser usadas:
%   \cite{Weisburd2022}, \cite{wang2019deep}, \cite{rotaru2022event},
%   \cite{albors_zumel_deep_2025}, \cite{papachristos_promises_2022}

\textbf{[SEÇÃO PENDENTE]} -- Redigir fundamentação sobre análise espacial do crime, representação espaço-temporal em grids, evolução de métodos de predição criminal (de modelos estatísticos a deep learning) e trabalhos relacionados que empregaram ConvLSTM ou arquiteturas similares para essa tarefa.


\subsection{Redes Neurais Artificiais}

As Redes Neurais Artificiais (RNAs) são sistemas computacionais baseados em modelos matemáticos inspirados no sistema nervoso, buscando aproximar, de forma abstrata, mecanismos de processamento e aprendizagem associados ao cérebro humano \cite{sala2019}. Um marco histórico ocorreu em 1943, quando Warren McCulloch e Walter Pitts propuseram um dos primeiros modelos formais de neurônio artificial através de um modelo matemático que estabeleceu as premissas fundamentais para futuras implementações computacionais \cite{mc1943}. Essas redes operam por processamento paralelo e distribuído, no qual múltiplas unidades simples atuam simultaneamente e de forma interconectada.

Em termos de funcionamento, uma RNA organiza essas unidades em camadas. Cada neurônio recebe valores de entrada, realiza uma combinação ponderada por pesos sinápticos, adiciona um viés e aplica uma função de ativação não linear, gerando um sinal que é propagado para os neurônios da camada seguinte. Essa composição em camadas permite o aprendizado de representações hierárquicas, onde camadas iniciais tendem a capturar padrões mais simples, enquanto camadas posteriores combinam esses padrões para modelar relações mais complexas \cite{gupta2018}.

Além da estrutura em camadas, o desempenho das RNAs depende do processo de treinamento, no qual os pesos e vieses são ajustados para minimizar o erro entre as saídas previstas e os valores reais. Em geral, esse ajuste é realizado por métodos de otimização baseados em gradiente, utilizando o algoritmo de retropropagação do erro (backpropagation) para calcular, de forma eficiente, como cada parâmetro da rede contribui para a função de perda \cite{rumelhart1986learning}. Os diferentes arranjos dessas camadas, bem como a forma como os parâmetros são conectados e compartilhados, dão origem a arquiteturas adequadas a diferentes tipos de dados e tarefas. Redes feedforward processam a informação em um fluxo unidirecional, sendo empregadas em problemas gerais de classificação e regressão. 

Quando os dados apresentam estrutura espacial, as Redes Neurais Convolucionais (CNNs) exploram a organização em grade por meio de operações de convolução e compartilhamento de pesos, tornando-se eficazes para identificar padrões locais e construir representações hierárquicas em imagens e mapas \cite{lecun2015}. Já nos cenários em que a informação é sequencial, as Redes Neurais Recorrentes (RNNs) incorporam um estado interno que é atualizado ao longo do tempo, permitindo a modelagem de dependências temporais entre observações. 
Contudo, RNNs tradicionais podem apresentar limitações para aprender dependências de longo prazo, motivando o desenvolvimento de arquiteturas com mecanismos de memória e portas de controle, como as Long Short-Term Memory (LSTM) \cite{lipton2015}.

%% ================================================================
%% TODO: Seção detalhada sobre LSTM
%% ================================================================
\subsection{TODO -- Long Short-Term Memory (LSTM)}
\label{Sec:todo_lstm}

% TODO: Esta subseção precisa ser escrita. A LSTM é base direta da ConvLSTM
% e merece tratamento próprio. Deve cobrir:
%
% 1. Problema do gradiente desvanecente (vanishing gradient) nas RNNs
%    tradicionais e por que a LSTM foi proposta para resolvê-lo.
%
% 2. Arquitetura da célula LSTM com suas 4 portas:
%    - Porta de esquecimento (forget gate): f_t = sigma(W_f * [h_{t-1}, x_t] + b_f)
%    - Porta de entrada (input gate): i_t = sigma(W_i * [h_{t-1}, x_t] + b_i)
%    - Candidata ao estado celular: g_t = tanh(W_g * [h_{t-1}, x_t] + b_g)
%    - Atualização do estado celular: c_t = f_t * c_{t-1} + i_t * g_t
%    - Porta de saída (output gate): o_t = sigma(W_o * [h_{t-1}, x_t] + b_o)
%    - Estado oculto: h_t = o_t * tanh(c_t)
%
% 3. Incluir diagrama da célula LSTM (recomendado).
%
% 4. Breve menção a variantes (GRU, Peephole LSTM) se desejar.
%
% Referências: \cite{lipton2015}, Hochreiter & Schmidhuber (1997) [adicionar à .bib]

\textbf{[SEÇÃO PENDENTE]} -- Detalhar a arquitetura LSTM: problema do gradiente desvanecente, mecanismo de portas (forget, input, output), equações do estado celular e do estado oculto. Incluir diagrama da célula LSTM.


Essas distinções tornam-se especialmente relevantes em problemas espaço-temporais, nos quais é necessário capturar simultaneamente padrões no espaço e dinâmicas de tempo. Nesse contexto, arquiteturas híbridas como a ConvLSTM combinam convoluções (para modelar as correlações espaciais em grids) com recorrência (para representar sequências de mapas), sendo, portanto, adequadas para tarefas de previsão em dados espaço-temporais \cite{shi2015}.


\subsection{ConvLSTM}

A ConvLSTM (Convolutional Long Short-Term Memory) é uma arquitetura derivada das redes LSTM, proposta por Shi et al. \cite{shi2015}, como já citado, para lidar com sequências espaço-temporais, ou seja, que possuem simultaneamente padrões espaciais e dinâmicas de tempo. Diferentemente da LSTM, em que as transformações entre entrada e estado são realizadas por operações totalmente conectadas, a ConvLSTM substitui essas operações por convoluções, preservando a estrutura espacial dos dados ao longo do tempo. Isso permite que o modelo capture tanto a evolução temporal quanto as correlações espaciais. 

As convoluções consistem na aplicação de pequenos filtros na entrada, também chamados de kernels, que analisam partes pequenas da grade de dados. Esses filtros identificam a concentração, mudanças e semelhanças entre a região analisada e as regiões vizinhas. À medida que percorrem a entrada, as características capturadas são organizadas em diferentes posições do espaço, formando representações que destacam onde determinados padrões ocorrem. 

%% ================================================================
%% TODO: Equações formais da ConvLSTM
%% ================================================================
% TODO: Adicionar as equações formais da ConvLSTM nesta subseção.
% A diferença para a LSTM é que as multiplicações matriciais (W * x)
% são substituídas por convoluções (W ** x):
%
%   i_t = sigma(W_{xi} ** X_t + W_{hi} ** H_{t-1} + b_i)
%   f_t = sigma(W_{xf} ** X_t + W_{hf} ** H_{t-1} + b_f)
%   C_t = f_t . C_{t-1} + i_t . tanh(W_{xc} ** X_t + W_{hc} ** H_{t-1} + b_c)
%   o_t = sigma(W_{xo} ** X_t + W_{ho} ** H_{t-1} + b_o)
%   H_t = o_t . tanh(C_t)
%
% onde ** denota convolução e . denota produto de Hadamard.
% Destacar que X_t, H_t, C_t são tensores 3D (canais x altura x largura),
% o que preserva a estrutura espacial.
% Incluir diagrama comparativo LSTM vs ConvLSTM (recomendado).
%
% Referência: \cite{shi2015}

\textbf{[PENDENTE]} -- Incluir equações formais da ConvLSTM (portas com convolução em vez de multiplicação matricial) e, idealmente, diagrama comparativo LSTM vs ConvLSTM.


%% ================================================================
%% TODO: Técnicas de regularização e treinamento
%% ================================================================
\subsection{TODO -- Técnicas de regularização e otimização}
\label{Sec:todo_regularizacao}

% TODO: Esta subseção precisa ser escrita. Conceitos usados no modelo
% mas que não foram fundamentados na revisão de literatura:
%
% 1. Dropout: desligamento aleatório de neurônios durante o treinamento
%    para prevenir sobreajuste (Srivastava et al., 2014) [adicionar à .bib].
%
% 2. Batch Normalization: normalização das ativações entre camadas
%    para estabilizar e acelerar o treinamento (Ioffe & Szegedy, 2015)
%    [adicionar à .bib].
%
% 3. Funções de ativação: Sigmoid, ReLU, Tanh -- usadas nas portas da
%    ConvLSTM e entre camadas. Breve descrição de cada uma.
%
% 4. Otimizador Adam: combinação de Momentum e RMSProp (Kingma & Ba, 2015)
%    [adicionar à .bib].
%
% 5. Learning rate scheduling: Cosine Annealing (Loshchilov & Hutter, 2017)
%    [adicionar à .bib].

\textbf{[SEÇÃO PENDENTE]} -- Fundamentar Dropout, Batch Normalization, funções de ativação (sigmoid, ReLU, tanh), otimizador Adam e escalonamento de taxa de aprendizado (Cosine Annealing). Todos são usados no modelo mas não estão no referencial.


%% ================================================================
%% TODO: Métricas de avaliação
%% ================================================================
\subsection{TODO -- Métricas de avaliação para classificação binária}
\label{Sec:todo_metricas}

% TODO: Esta subseção precisa ser escrita. As métricas são usadas
% extensivamente na avaliação mas nunca definidas formalmente:
%
% 1. Matriz de confusão: TP, FP, TN, FN.
%
% 2. Precision = TP / (TP + FP)
%
% 3. Recall (Sensibilidade) = TP / (TP + FN)
%
% 4. F1-Score = 2 * (Precision * Recall) / (Precision + Recall)
%
% 5. Threshold de classificação: como logits/probabilidades são
%    convertidos em previsão binária via limiar.
%
% 6. Desbalanceamento de classes: definição formal, impacto nas métricas
%    (accuracy paradox), e estratégias de mitigação:
%    - Ponderação da função de perda (pos_weight)
%    - Focal Loss (Lin et al., 2017) [adicionar à .bib]
%    - Oversampling/undersampling
%
% 7. BCEWithLogitsLoss: combinação de sigmoid + binary cross-entropy
%    com estabilidade numérica.

\textbf{[SEÇÃO PENDENTE]} -- Definir formalmente matriz de confusão, precision, recall, F1-score, threshold de classificação, desbalanceamento de classes e estratégias de mitigação (pos\_weight, focal loss). Também definir BCEWithLogitsLoss.


%% ================================================================
%% TODO: Trabalhos relacionados
%% ================================================================
\subsection{TODO -- Trabalhos relacionados}
\label{Sec:todo_trabalhos_relacionados}

% TODO: Esta subseção precisa ser escrita. Deve discutir os trabalhos
% que servem de base e comparação para este estudo:
%
% 1. Albors Zumel, Tizzoni & Campedelli (2025) -- artigo base deste TCC:
%    - ConvLSTM para crime forecasting em 4 cidades dos EUA
%    - Grid de 0.2 km², previsão 12h à frente, 14 dias de lookback
%    - Comparação com baselines (logistic regression, random forest, LSTM)
%    - Resultados: DL supera baselines, mas precisão geral ainda limitada
%    - Contribuição de mobilidade + sociodemografia
%    \cite{albors_zumel_deep_2025}
%
% 2. Wang et al. (2019) -- deep learning para crime forecasting em tempo real:
%    - Formulação ternária do problema
%    \cite{wang2019deep}
%
% 3. Rotaru et al. (2022) -- predição event-level com viés de enforcement:
%    - Análise de viés policial nos dados e nos modelos
%    \cite{rotaru2022event}
%
% 4. Papachristos (2022) -- promessas e perigos da predição de crime:
%    - Aspectos éticos e limitações
%    \cite{papachristos_promises_2022}
%
% Comparar: dados usados, granularidade, arquitetura, métricas, resultados.

\textbf{[SEÇÃO PENDENTE]} -- Revisar e comparar os trabalhos de referência: Albors Zumel et al. (2025) como base metodológica, Wang et al. (2019) para formulação do problema, Rotaru et al. (2022) para questões de viés, e Papachristos (2022) para aspectos éticos. Comparar dados, granularidade, arquiteturas e resultados obtidos.


%% ================================================================
%%  METODOLOGIA
%% ================================================================
\section{Metodologia}
\label{Sec:metodologia}

Esta seção detalha as etapas do pipeline experimental implementado neste trabalho, desde a obtenção dos dados brutos até a avaliação do modelo treinado. Todas as etapas foram desenvolvidas em Python, utilizando Jupyter Notebooks e bibliotecas como pandas, geopandas, osmnx, shapely, folium, PyTorch e scikit-learn.

\subsection{Dados de origem}
\label{Sec:dados}

Os dados de ocorrências criminais foram obtidos da Polícia Militar de Alagoas por meio de parceria institucional, mediante assinatura de termo de confidencialidade. Em razão dessa restrição, a base de dados original não pode ser disponibilizada publicamente. Os dados correspondem a registros georreferenciados (latitude e longitude) de crimes violentos contra o patrimônio nas cidades de Maceió e Arapiraca, cobrindo o período de 2012 a 2022. O conjunto original contém 137.300 registros com 30 variáveis, incluindo natureza do fato, data/hora, coordenadas geográficas e cidade.

Embora os dados brutos não possam ser divulgados, todo o código-fonte utilizado para o pré-processamento, a construção dos grids, a preparação das sequências e o treinamento do modelo está disponível publicamente em repositório no GitHub\footnote{\url{https://github.com/TODO-inserir-link-do-repositorio}}, permitindo a reprodutibilidade metodológica do trabalho.

\subsection{Pré-processamento dos dados}
\label{Sec:preprocessamento}

O pré-processamento foi realizado em múltiplas etapas encadeadas:

\subsubsection{Carregamento e limpeza inicial}

Os arquivos CSV de entrada foram carregados e concatenados. As colunas de data foram convertidas para o formato datetime (formato \texttt{dd/mm/yyyy HH:MM:SS}) e as coordenadas de latitude e longitude, originalmente com separador decimal em vírgula, foram convertidas para tipo numérico. Registros com data ou coordenadas inválidas (valores ausentes) foram removidos.

\subsubsection{Filtragem por tipo de crime e agrupamento}

Das 23 naturezas de fato distintas presentes nos dados, foram selecionadas 10 categorias relevantes de roubo, que foram agrupadas em 5 macro-categorias:
\begin{itemize}
  \item \textbf{street\_robbery}: roubo a transeunte;
  \item \textbf{vehicle\_robbery}: roubo de veículo (moto, passeio e outros);
  \item \textbf{public\_transport\_robbery}: roubo a transporte coletivo (urbano e rodoviário);
  \item \textbf{commercial\_robbery}: roubo a casa comercial, correspondente bancário e correios;
  \item \textbf{residential\_robbery}: roubo a residência.
\end{itemize}
Após essa filtragem e agrupamento, restaram 119.455 registros, sendo 72.787 em Maceió e 17.421 em Arapiraca, que foram separados em arquivos distintos por cidade.

\subsubsection{Filtragem espacial com polígono municipal}

Para cada cidade, o polígono do limite municipal foi obtido via OpenStreetMap (usando a biblioteca osmnx). Em seguida, realizou-se um spatial join entre os pontos de ocorrência e o polígono, mantendo apenas os registros cujas coordenadas caem efetivamente dentro do município. Para Maceió, 83.512 de 83.612 registros foram mantidos (100 descartados por estarem fora do polígono); para Arapiraca, 19.141 de 19.197 (56 descartados). Os resultados foram visualizados em mapas interativos (folium) para verificação.

\subsection{Construção da representação em grids}
\label{Sec:grids}

\subsubsection{Geração da grade espacial}

Para Maceió, uma grade regular de $39 \times 39$ células foi sobreposta ao bounding box do polígono municipal, resultando em 1.521 células. Cada ocorrência criminal foi atribuída à célula correspondente por meio de verificação de contenção geométrica (point-in-polygon).

% TODO: Calcular e informar a área aproximada de cada célula em km².
% O bounding box de Maceió tem extensão aprox. de ~0.25° lat × ~0.15° lon.
% Com 39 divisões, cada célula tem lado aprox. de X km, resultando em
% uma área de ~Y km². Comparar com o estudo de referência (Albors Zumel
% et al., 2025), que utilizou células de 0,2 km² (0,077 sq. miles).
% Justificar a escolha de 39×39 (ou indicar que foi experimental).
\textbf{[PENDENTE]} -- Informar a área aproximada de cada célula (em km\textsuperscript{2}) e justificar a escolha do tamanho $39 \times 39$. Comparar com o grid de 0,2~km\textsuperscript{2} usado por Albors Zumel et al. (2025).

\subsubsection{Máscara de validade}

Uma máscara binária de $39 \times 39$ foi gerada, indicando quais células do grid possuem interseção com o polígono municipal. Para Maceió, 811 das 1.521 células foram marcadas como válidas. Essa máscara é utilizada na etapa de avaliação para ignorar células que estão fora dos limites da cidade.

\subsubsection{Agregação temporal}

As ocorrências foram contabilizadas por célula e por hora do ano, formando uma matriz esparsa \textit{célula} $\times$ \textit{hora}, para cada tipo de crime e para cada ano. Em seguida, essas matrizes horárias foram agregadas em janelas diárias (granularidade de 24h), resultando em sequências temporais com aproximadamente 365 passos por ano. As contagens dos 5 tipos de crime foram somadas e binarizadas (presença/ausência de crime: $\geq 1 \rightarrow 1$, caso contrário $\rightarrow 0$), produzindo um tensor final de dimensão $4.018 \times 39 \times 39 \times 1$ (dias $\times$ linhas $\times$ colunas $\times$ canais) para Maceió, abrangendo 2012--2022.

\subsection{Preparação para treinamento}
\label{Sec:preparacao}

\subsubsection{Janela deslizante e sub-grids}

Uma janela deslizante de 8 passos (7 dias de lookback + 1 dia de alvo) foi aplicada sobre a sequência temporal. Para cada amostra, em vez de utilizar o grid completo de $39 \times 39$, foram extraídos 5 sub-grids de $16 \times 16$ em posições aleatórias dentro da grade, ampliando o volume de amostras de treinamento e permitindo que o modelo opere sobre regiões locais da cidade.

% TODO: Justificar as seguintes escolhas de hiperparâmetros do pipeline:
% - Por que lookback de 7 dias (e não 14, como no artigo de referência)?
% - Por que sub-grids de 16×16? Qual a relação com o grid total de 39×39?
% - Por que 5 sub-grids por amostra temporal? Isso multiplica os dados por 5.
% - Essas escolhas foram baseadas no artigo de referência, em experimentos
%   preliminares, ou em limitações computacionais?
\textbf{[PENDENTE]} -- Justificar a escolha do lookback de 7 dias, o tamanho do sub-grid ($16 \times 16$) e a quantidade de 5 sub-grids por amostra. Indicar se foram baseados na literatura, em experimentos ou em restrições computacionais.

\subsubsection{Divisão cronológica treino/validação}

A divisão dos dados respeita rigorosamente a ordem temporal para evitar vazamento de informação. Os primeiros 90\% da sequência são utilizados para treinamento e os 10\% finais para validação. Para Maceió, isso resultou em 17.985 amostras de treinamento e 1.925 amostras de validação (cada amostra com formato $7 \times 16 \times 16 \times 1$ para a entrada e $16 \times 16 \times 1$ para o alvo). Os dados foram salvos no formato \texttt{.npz} (arrays comprimidos NumPy). Valores ausentes (NaN) remanescentes foram substituídos por zero antes do treinamento.

% TODO: Discutir as seguintes questões sobre a divisão dos dados:
% - A divisão é apenas treino/validação (90/10). Não há conjunto de teste
%   independente. Isso significa que os resultados reportados são sobre
%   validação, não sobre dados completamente inéditos. Considerar se é
%   necessário adicionar um split de teste (ex.: 80/10/10) ou justificar
%   por que a validação cronológica é suficiente.
% - O split 90/10 coloca ~10 anos em treino e ~1 ano em validação.
%   Indicar quais anos aproximados compõem cada partição.
\textbf{[PENDENTE]} -- Discutir a ausência de conjunto de teste separado (apenas treino/validação). Justificar ou propor um split treino/validação/teste. Indicar quais anos compõem cada partição.


%% ================================================================
%%  MODELO E TREINAMENTO
%% ================================================================
\section{Arquitetura do modelo e treinamento}
\label{Sec:modelo}

\subsection{Arquitetura ConvLSTM implementada}

O modelo foi implementado em PyTorch e consiste em três componentes principais:

\begin{enumerate}
  \item \textbf{ConvLSTM empilhada}: 3 camadas ConvLSTM com 7 canais ocultos cada e kernels de $3 \times 3$. Cada camada ConvLSTM é seguida por ativação ReLU e normalização por lote (BatchNorm3d).
  \item \textbf{Dropout}: camada de dropout com taxa $p = 0{,}8$ aplicada sobre o estado oculto da última camada.
  \item \textbf{Convolução final}: uma convolução $3 \times 3$ com 1 canal de saída (padding ``same''), produzindo o mapa de previsão para o dia seguinte.
\end{enumerate}

O modelo possui 9.262 parâmetros treináveis. A entrada tem formato $(B, 7, 16, 16, 1)$ (batch, sequência, altura, largura, canais) e a saída tem formato $(B, 16, 16)$, representando as probabilidades de ocorrência de crime para cada célula no próximo passo temporal.

% TODO: Incluir justificativa dos hiperparâmetros da arquitetura:
% - Por que 3 camadas ConvLSTM? (profundidade)
% - Por que 7 canais ocultos? (capacidade representacional)
% - Por que kernel 3×3?
% - Por que dropout de 0.8? (valor muito alto, justificar)
% - Essas escolhas foram tuned? Grid search? Baseadas no artigo de referência?
% - Incluir diagrama/figura da arquitetura do modelo (recomendado para TCC).
\textbf{[PENDENTE]} -- Justificar a escolha dos hiperparâmetros da arquitetura (3 camadas, 7 canais, kernel 3$\times$3, dropout 0,8). Indicar se foram baseados no artigo de referência ou em tuning. Incluir diagrama/figura da arquitetura (recomendado).

\subsection{Função de perda e otimização}

Dada a severa esparsidade dos dados ($\sim$2\% das células apresentam crime em qualquer dia), foi utilizada a função de perda \texttt{BCEWithLogitsLoss} com \texttt{pos\_weight} calculado como a razão entre o número de amostras negativas e positivas no conjunto de treinamento ($\approx$48,2 para Maceió). Essa ponderação atribui maior importância aos casos positivos durante o treinamento.

O otimizador utilizado foi o Adam com taxa de aprendizado de $10^{-5}$, e um escalonador do tipo Cosine Annealing com $T_{\max} = 20$ épocas. O treinamento foi conduzido por 200 épocas com batch size de 55, utilizando GPU (NVIDIA GeForce RTX 4070 Laptop GPU) e semente aleatória fixa (42) para reprodutibilidade.

\subsection{Avaliação com tolerância espacial}
\label{Sec:avaliacao}

Além das métricas tradicionais (precision, recall, F1), foi implementada uma avaliação com tolerância espacial, conforme proposto por \citeasnoun{albors_zumel_deep_2025}. Nessa variante, um falso positivo cuja célula vizinha (distância de Chebyshev $\leq 1$) contenha um crime real é reclassificado como acerto espacial (``true positive espacial''). Isso resulta em métricas \textit{precision\_spatial}, \textit{recall\_spatial} e \textit{f1\_spatial}, que reconhecem previsões geograficamente próximas ao evento real.

A avaliação foi realizada sobre o conjunto de validação, varrendo thresholds de classificação de 0,50 a 1,00 (incrementos de 0,01). A máscara de validade municipal foi aplicada para que apenas as células dentro dos limites da cidade contribuíssem para as métricas.



